\subsection{A prescribed-point criterion from an automorphism orbit}

\begin{claim}{B.145}{Proved}{Long automorphism orbits force pointed class-number witnesses}
Let $C/k$ be smooth, proper and geometrically connected of genus
$g\ge2$ over $k=\mathbf F_q$ of characteristic $p$.
Let $P\in C(k)$, let $\ell\ne p$ be an odd prime, and let
$\chi:J_C(k)\twoheadrightarrow\mathbf Z/\ell$.
Suppose there is $\tau\in\operatorname{Aut}_k(C)$ and
$u\in(\mathbf Z/\ell)^\times$ such that
\[
 \chi\circ\tau_*=u\chi,\qquad
 u\in\langle p\rangle,\qquad
 m:=\operatorname{ord}_{\ell}(u)>g-1 .
                                                            \tag{B.145.1}
\]
Put $a=\chi([\tau P-P])$.
Index the constant twists of the associated pointed Lang cover by
$t\in\mathbf Z/\ell$, so that a closed point of degree $d$
acquires the added Frobenius translation $dt$.
The only possibly bad twist, meaning one whose relative class
number is divisible by $p$, is
\[
                              t=\frac{a}{u-1}.
                                                            \tag{B.145.2}
\]
In particular, if $a\ne0$, the specified $P$-split cover, at
$t=0$, is a class-number witness already over $k$.
That same cover remains a witness over every constant extension
of degree prime to $\ell$.
There is no restriction that the order of $\tau$ be prime to $p$.

An explicit genus-three plane-quartic example is
\[
 C:\ Y^3Z-YZ^3=X^4,\qquad
 k=\mathbf F_{27}=\mathbf F_3[t]/(t^3-t-1),\qquad
 P=(x,y)=(1,t),\qquad \ell=7 .
                                                            \tag{B.145.3}
\]
The automorphism is $\tau(x,y)=(x,y+1)$.
The character specified below has $u=2=3^2\bmod7$ and $a=1$.
Its constant cyclic cover has genus fifteen and
\[
\begin{split}
 \#J_C(k)&=18928,\\
 \frac{\#J_D(k)}{\#J_C(k)}&=427439613252172921
                                  \equiv1\pmod3,\\
 \#J_D(k)&=8090576999637129048688 .
\end{split}                                                   \tag{B.145.4}
\]
At this particular field and marking the least admissible
witness prime is seven.
All twenty-four points of $C(k)\setminus C(\mathbf F_3)$
are certified by this character construction.
The four original $\mathbf F_3$ pointings are not certified
by the nonzero-shift hypothesis, but the separate finite twist
calculation also gives them a degree-seven witness over $k$,
with relative class number $121740744925904896$.
For each original pointing the first field admitting an
odd-prime class-number witness has degree three over
$\mathbf F_3$, and seven is the least prime at that field.
For the displayed seven twists over $k$, base change to
$\mathbf F_{27^n}$ has relative class number prime to three
exactly when $4\nmid n$ for twist one, and exactly when
$28\nmid n$ for every other twist.
At the stated field, the split pointed models are twist one
for the four original markings and twists zero, four and six
for the twenty-four new markings.  Twists two, three and five
have no split $k$-rational marking; their assertion is only
the relative class-number unit test.
\end{claim}
\begin{proof}
\medskip\noindent\textbf{The general character polynomial and its degree modulo $p$.}
We use the connected constant Lang cover attached to $\chi$,
with its chosen rational point over $P$, as constructed in B.2.
Fix a primitive $\ell$-th root of unity $\zeta$.
For $i\ne0\bmod\ell$ let
\[
 A_d(i)=\sum_{\substack{D\ge0\\ \deg D=d}}
              \zeta^{\,i\chi([D-dP])},\qquad
                 L_i(T)=\sum_{d\ge0}A_d(i)T^d .
                                                            \tag{B.145.5}
\]
The sum is over effective divisors defined over $k$.
Its Euler product is the unramified character factor of the cover.
We recall the exact Riemann--Roch reduction needed here.

Identify $\operatorname{Pic}^d(C)(k)$ with $J_C(k)$ using $dP$.
The number of effective divisors in a line-bundle class $M$
is $(q^{h^0(M)}-1)/(q-1)$.
For $d>2g-2$, that number is the same in every class, so
nontrivial character orthogonality gives $A_d(i)=0$.
For $g\le d\le2g-2$, Riemann--Roch and the same orthogonality
give the exact identity
\[
 A_d(i)=q^{d+1-g}\zeta^{\,i c}
                      A_{2g-2-d}(-i),\qquad
             c=\chi([K_C-(2g-2)P]).
                                                            \tag{B.145.6}
\]
Indeed replace $M$ by $K_C\otimes M^{-1}$.
Its section dimension differs from $h^0(M)$ by $g-1-d$.
The contribution of the constant term in the numerator
$q^{h^0(M)}-1$ vanishes after summing the character.
Factoring out $q^{d+1-g}$ leaves the effective-divisor count
for $K_C\otimes M^{-1}$, with character multiplier
$\zeta^{ic}$.
The rational point ensures these line-bundle classes and their
linear systems are defined over $k$.
Consequently
\[
                   \deg(\overline L_i)\le g-1,\qquad
                            \overline L_i(0)=1
                                                            \tag{B.145.7}
\]
after reduction at any prime above $p$.
The polynomial is therefore not identically zero.

\medskip\noindent\textbf{The affine action on twists.}
The automorphism $\tau$ permutes the effective divisors of
each degree.  For any such divisor,
\[
 [\tau D-dP]=\tau_*[D-dP]+d[\tau P-P].
\]
Applying $\chi$ and summing gives
\[
                  A_d(iu)=\zeta^{-ida}A_d(i),\qquad
                         L_{iu}(T)=L_i(\zeta^{-ia}T).
                                                            \tag{B.145.8}
\]
Choose $j\ge1$ with $p^j=u\bmod\ell$.
If the $i$-th character factor of twist $t$ vanishes modulo $p$,
then $\overline L_i(\bar\zeta^{it})=0$.
Raising this equality to the $p^j$-th power and using
(B.145.8) gives
\[
              \overline L_i(\bar\zeta^{\,i(ut-a)})=0.
                                                            \tag{B.145.9}
\]
Here coefficients are sums of rational-integer multiples of
roots of unity, so the same power changes the character index
from $i$ to $iu$.
The image $\bar\zeta$ has exact order $\ell$, because
$\ell\ne p$.

The affine map $t\mapsto ut-a$ has the unique fixed point
$a/(u-1)$, and every other orbit has length $m$.
Indeed, after translating by that fixed point, it is
multiplication by $u$ on $\mathbf F_\ell$.
The assumption $m>g-1\ge1$ ensures $u\ne1$.
For any nonfixed $t$, iteration of (B.145.9) would give
$m$ distinct roots of the nonzero polynomial
$\overline L_i$, contradicting (B.145.7).
This excludes a vanishing factor for every primitive index $i$,
with the same unique exceptional twist for all of them.

As in the direct Euler-factor calculation of B.141, the product
of the nontrivial character factors at one is exactly
$\#J_{D_t}(k)/\#J_C(k)$.
It is an integer, and it is divisible by $p$ precisely when one
of those reduced factors vanishes.
This proves (B.145.2), including the assertion for the
prescribed $P$-split twist.

For a degree-$n$ constant extension with $\ell\nmid n$,
the same cover has character
$\chi_n=\chi\circ\operatorname{Nm}_{k_n/k}$.
Norm commutes with $\tau_*$, so the eigencharacter relation
still holds with $u$.
Moreover
$\chi_n([\tau P-P])=n\chi([\tau P-P])=na\ne0$.
The character is nonzero since its restriction to $J_C(k)$
is $n\chi$.
The already proved argument over $k_n$ proves persistence.

\medskip\noindent\textbf{A nonempty genus-three example.}
The three partial derivatives of
$Y^3Z-YZ^3-X^4$ in characteristic three are
$-X^3,-Z^3,Y^3$; they have no common projective zero.
Thus (B.145.3) is a smooth plane quartic of genus three.
Its unique point at infinity is $O=(0:1:0)$.
The line $Z=0$ cuts out $4O$, so $K_C\sim4O$.
There is a degree-two morphism
\[
 \pi:C\longrightarrow E,\qquad
 E:\ v^2=u^3-u,\qquad (u,v)=(y,x^2),
                                                            \tag{B.145.10}
\]
which sends $O$ to the elliptic origin and intertwines $\tau$
with $\tau_E(u,v)=(u+1,v)$.

The point $R=(t,1)$ belongs to $E(k)$.
The doubling law gives
$2R=(t+1,1)$, $4R=(t+2,1)$ and $8R=R$.
Hence $R$ has exact order seven.
Over $\mathbf F_3$, the elliptic curve has four points and
Frobenius trace zero.  Its two Frobenius roots have sum zero,
so their cubes also have sum zero and $\#E(k)=28$.
All four points of $E[2]$ are already rational over
$\mathbf F_3$.
It follows that
\[
                    E(k)=E[2]\oplus\langle R\rangle .
\]
Define $\chi_E:E(k)\to\mathbf Z/7$ by
$\chi_E(R)=1$ and $\chi_E(E[2])=0$.
Since $\tau_E(R)=2R$ and $\tau_E$ permutes $E[2]$,
$\chi_E\circ\tau_E=2\chi_E$.
Set $\chi=\chi_E\circ\pi_*:J_C(k)\to\mathbf Z/7$.
For $P=(1,t)$, one has $\pi(P)=R$ and
\[
                 \chi([\tau P-P])=\chi_E(2R-R)=1.
\]
Thus $\chi$ is surjective, its eigenvalue is $u=2=3^2\bmod7$,
and $m=3>2=g-1$.
The cover is a witness by the general theorem.
This verifies the criterion's hypotheses geometrically and
arithmetically, before any finite class-number computation.

\medskip\noindent\textbf{The complete finite arithmetic certificate.}
The program \path{finite_checks/code/b145_quartic_example.py}
uses the displayed cubic field and
$\mathbf F_{729}=k[s]/(s^2+1)$.
Its elements are encoded by $a+27b$ with $a,b\in k$;
the $27$-power Frobenius sends $s$ to $-s$.
It retains every degree-one point and every degree-two
closed-point orbit used below, together with their elliptic
norms and character values.
The elliptic addition and character tables are also retained
and checked on every pair.

The exact curve counts over $\mathbf F_3,\mathbf F_9,k$ and
$\mathbf F_{729}$ are respectively $4,4,28,676$.
For the first three, Newton's identities give the
$\mathbf F_3$ Frobenius polynomial
\[
                         (X^2-3)^2(X^2+3).
\]
Cubing its roots gives $(X^2-27)^2(X^2+27)$ over $k$,
and hence $\#J_C(k)=(-26)^2\cdot28=18928$.
For the character based at $P$, the degree-one and
degree-two effective-divisor occupancy vectors are
\[
\begin{split}
 (n_0,\ldots,n_6)&=(8,8,0,8,0,0,4),\\
 (m_0,\ldots,m_6)&=(116,116,116,116,116,34,116).
\end{split}                                                   \tag{B.145.11}
\]
The second vector includes the $28\cdot29/2=406$ unordered
pairs of rational points, with repetition, and all
$(676-28)/2=324$ closed points of degree two.
The character of a degree-two closed point is computed by
the elliptic norm of its two geometric points minus $2R$.
Thus these are the complete coefficient sums in (B.145.5),
not point-count proxies for the degree-two coefficient.

One has $c=\chi([4O-4P])=-4=3\bmod7$.
Writing
\[
       S=8+8\zeta+8\zeta^3+4\zeta^6,
\]
the resulting full character factor is
\[
 L_1(T)=1+ST-82\zeta^5T^2+
             27\zeta^3 S(\zeta^{-1})T^3+729\zeta^3T^4 .
                                                            \tag{B.145.12}
\]
Here the degree-two vector reduces to $-82\zeta^5$
using $1+\zeta+\cdots+\zeta^6=0$, and the two high-degree
coefficients follow from (B.145.6).
At one, it is represented by
\[
 A(Z)=225+8Z+216Z^2+953Z^3+108Z^4-82Z^5+4Z^6 .
\]
The complete integer norm is
\[
 \operatorname{Res}_Z(1+Z+\cdots+Z^6,A(Z))
                         =427439613252172921 .
                                                            \tag{B.145.13}
\]
The program computes this norm both as an integer resultant
and as the determinant of multiplication on
$\mathbf Z[Z]/\Phi_7$; it also retains the full relative
zeta numerator of degree twenty-four.
As a separate unit check, $\operatorname{ord}_7(3)=6$,
so $\Phi_7$ is irreducible modulo three, whereas
$A(Z)\bmod3=2Z+2Z^3+2Z^5+Z^6$ is nonzero and is not
a multiple of $\Phi_7$.
This proves (B.145.4).
Since $5\nmid18928$ and five is the only smaller odd prime
different from three, no admissible smaller prime can give
a constant pointed cyclic quotient at this field.
Seven is therefore the least witness prime at this field.

For clarity about the marking, $C(\mathbf F_3)$ consists of
$O$ and $(0,j)$ for $j\in\mathbf F_3$.
Their images under $\pi$ lie in $E[2]$, so their
$\chi([\tau P-P])$ values are zero.
The complete $k$-point table gives nonzero
$\chi_E(\pi(P))$ at every other point: there are eight
points at each of the three values $1,2,4$.
These are exactly the other twenty-four points.
For each, the automorphism shift is
$(2-1)\chi_E(\pi(P))\ne0$, proving the assertion for
all of those new pointings.
The three sets are permuted by $\tau$, and the pointed
character factors depend only on the displayed value,
so their relative class number is the same integer.

\medskip\noindent\textbf{The remaining twist and the original pointings.}
The criterion leaves twist one undecided in this example,
since $a/(u-1)=1$.
The complete finite calculation settles it too.
The exact norms for all twists are
\[
\begin{array}{c|r}
 t & \#J_{D_t}(k)/\#J_C(k)\\ \hline
 0,4,6&427439613252172921\\
 1&121740744925904896\\
 2,3,5&56515205339274361 .
\end{array}                                                  \tag{B.145.14}
\]
Every displayed integer is a three-adic unit.
For an original $\mathbf F_3$ point $P_0$, one has
$\chi_E(\pi(P_0))=0$, whereas $\chi_E(\pi(P))=1$.
Changing the pointing from $P$ to $P_0$ therefore adds one
to every degree-one exponent and $d$ to every degree-$d$
exponent.  Its split model is exactly twist one.
This proves the witness assertion for all four original
pointings, by the finite calculation rather than the
nonzero-shift criterion.

There is also an exact test through every constant extension
for each of these seven covers.
Work now in characteristic three and put
\[
 U=\bar\zeta+\bar\zeta^2+\bar\zeta^4,\qquad
                    \eta=1+2U .
\]
The quadratic-period identity $U^2+U+2=0$ gives
$\eta^2=-1$.
Reducing (B.145.12) and substituting $\bar\zeta T$ gives
\[
 \overline L_{1,1}(T)=1+\eta T-T^2=(1-\eta T)^2,
 \qquad
 \overline L_{1,t}(T)
                  =(1-\eta\bar\zeta^{\,t-1}T)^2 .
                                                            \tag{B.145.15}
\]
All other primitive character factors are its conjugates.
The nonzero reduced Frobenius roots consequently have
order four for $t=1$, and order twenty-eight for $t\ne1$,
since $\eta$ and a nontrivial seventh root have coprime
orders four and seven.
Under a degree-$n$ extension, every Frobenius root is
replaced by its $n$-th power.
Roots that reduce to zero cannot contribute a factor
divisible by three in the evaluation at one.
For the other roots, such divisibility occurs exactly
when their $n$-th power is one.
Hence
\[
 3\nmid\frac{\#J_{D_t}(\mathbf F_{27^n})}
                 {\#J_C(\mathbf F_{27^n})}
 \quad\Longleftrightarrow\quad
 \begin{cases}
  4\nmid n,&t=1,\\
  28\nmid n,&t\ne1 .
 \end{cases}                                                \tag{B.145.16}
\]
This is a proof for all extension degrees, not extrapolation
from a finite list of extensions.
The program separately checks the coefficient identities
in (B.145.15) and the exact orders four and twenty-eight
in $\mathbf F_3[Z]/\Phi_7$.

Finally the base Frobenius polynomial already computed gives
\[
 \#J_C(\mathbf F_3)=16,\qquad
 \#J_C(\mathbf F_9)=256,\qquad
 \#J_C(\mathbf F_{27})=16\cdot7\cdot13^2 .
                                                            \tag{B.145.17}
\]
A constant cyclic pointed cover of odd prime degree requires
a quotient of the rational Jacobian of that order.
The first two orders exclude all such covers.
At degree three the twist-one witness exists and the third
order makes seven the least admissible prime there.
This proves the claimed field-first minimum for every
original $\mathbf F_3$ marking.
\end{proof}

\paragraph{Exact scope and predecessor comparison.}
The general statement proves B1 for the entire marked-curve
family satisfying (B.145.1) with nonzero automorphism shift.
It includes B.141 by taking the hyperelliptic involution in
genus two, and the genus-three example proves that its scope
is larger.  It imposes a real automorphism and eigencharacter
condition; it does not manufacture such an automorphism for
an arbitrary curve, and it does not settle the unrestricted
threshold $p>(2g-2)^2$.
The finite quartic certificate treats new
$\mathbf F_{27}$ pointings and, by the separate twist calculation,
the first-field minimum for all four original $\mathbf F_3$
pointings of this quartic.  It gives the exact constant-extension
unit test for each of the seven displayed covers.
It does not determine the least prime when arbitrary extension
degrees are allowed, the least witness among all covers over
every extension, or the full B6 census.

\paragraph{Check record.}
Written and checked by GPT-6 Astra Pro, finite-checks reader,
18 November 2026; the separate effort setting is not exposed.
All finite inputs and complete outputs are retained with the
program and JSON certificate.
The infinite family is proved by the automorphism and
Riemann--Roch argument, not inferred from the finite example.
Independent second-reading events are recorded separately.
